\chapter{Referencial Teórico}
\label{cap:Teorico}


\section{Processamento de Linguagem Natural (PLN)}
\label{sec:Teorico.1}

O \ac{PLN} -- do inglês \textit{\ac{NLP}} -- pode ser definido como o processamento automático ou semi-automático da linguagem humana por sistemas computacionais. Trata-se de um campo multidisciplinar situado na interseção entre a ciência da computação, a \ac{IA} e a linguística, cujo objetivo é possibilitar a compreensão, interpretação e manipulação da linguagem humana, abrangendo, também, a geração de conteúdo que siga os moldes das comunicações interpessoais, servindo como ferramenta útil em interações significativas entre indivíduos e máquinas. \cite{PNL-Cambridge:}

O \ac{PLN} é intrinsecamente multidisciplinar, apropriando-se de conceitos e técnicas de diversas áreas para sua fundamentação teórica e aplicações práticas. Variando consideravelmente entre projetos e subáreas, é notável a influência de campos como o da teoria linguística, das ciências cognitivas e da psicologia, os quais contribuem para a compreensão do processamento da linguagem e da cognição humanas, oferecendo modelos que dão base aos algoritmos e técnicas de \ac{PLN}. Também é relevante o impacto do raciocínio lógico e das estruturas formais -- fundamentados na filosofia e na matemática -- os quais desempenham um papel crucial em áreas como a teoria formal das linguagens e a prova de teoremas.

Na ciência da computação, o \ac{PLN} tem suas raízes históricas na teoria das linguagens formais, sendo utilizada em campos como a concepção de compiladores e a interação humano-computador, baseando-se em regras formais e estatística. Todavia, com os avanços relativamente recentes em áreas como aprendizagem de máquina e redes neurais, bem como o aumento do poder de processamento dos computadores e do volume de dados disponíveis para análise, tornou-se possível o aprendizado automatizado de padrões que potencializou as técnicas de \ac{PLN}, de forma que este tem estado cada vez mais indissociável da \ac{IA}, tanto influenciando quanto sendo influenciado por esta.

\section{Representação Textual}
\label{sec:Teorico.2}
Sistemas computacionais operam, por natureza, exclusivamente sobre dados numéricos. Nesse contexto, a conversão de texto em representações numéricas torna-se um elemento indispensável para a realização de tarefas de \ac{PLN}. O texto, em sua forma bruta, constitui um construto humano altamente abstrato, desprovido de estrutura ou significado intrínsecos, inviabilizando o seu processamento direto, sem intervenções, por máquinas. Por essa razão, a transformação de informações textuais em formatos numéricos é um passo fundamental para permitir que sistemas computacionais identifiquem e interpretem padrões linguísticos, configurando-se como requisito essencial para uma ampla gama de aplicações no campo do \ac{PLN}. \cite{Representacao-Texto:1}

Esse processo, comumente denominado representação textual, desempenha um papel central, por exemplo, em atividades como a seleção de características para modelos e algoritmos de aprendizado de máquina. Ao converter textos em atributos numéricos compatíveis com as estruturas desses modelos, a representação textual viabiliza o processamento eficiente e a extração de padrões linguísticos relevantes.

De maneira geral, as abordagens para representação textual podem ser classificadas em duas categorias principais: representações discretas e representações distribuídas \cite{Representacao-Texto:1} -- (também referidas como ``representações locais'' \cite{Representacao-Texto:4} e ``representações contínuas'' \cite{Representacao-Discreta:1,Representacao-Distribuida:3}, respectivamente.

A representação textual discreta refere-se à transformação numérica do texto de forma que cada elemento textual (caracteres, sequências de caracteres, palavras isoladas, combinações de palavras, etc.) tenha uma correspondência clara e distinta com sua forma numérica, de modo que a codificação de um elemento seja completamente independente de outros. Essa abordagem é intuitiva e natural para representar o texto, uma vez que a língua e outras formas simbólicas possuem uma estrutura composicional, sendo um exemplo de uso de codificação simbólica discreta nas interações humanas, sejam elas interpessoais ou mediadas por máquinas \cite{Representacao-Discreta:1}.

Entre os aspectos positivos desse modelo de codificação, \citeonline{Representacao-Discreta:1} apontam que:
\begin{citacao}
    \textit{Natural language text is composed of strings of word tokens. These word tokens enable fast and exact processing. It is, for instance, straightforward to represent tokens in a computer by their indices. (...) Moreover, discrete representations have the advantage of being readily interpretable.} \footnote{Textos em linguagem natural são compostos por sequências de tokens de palavras. Esses tokens permitem um processamento rápido e preciso. Por exemplo, é relativamente simples representar tokens em um computador por meio de seus índices. (...) Além disso, representações discretas têm a vantagem de serem facilmente interpretáveis. (Tradução nossa)}
\end{citacao}


Tal abordagem é, de fato, extremamente simples em termos de implementação e interpretabilidade das representações; todavia -- dado que palavras são tratadas como símbolos atômicos e representadas também atomicamente, de modo que a ligação semântica entre elas não é capturada -- é ineficaz em capturar relações entre os elementos textuais representados, bem como em se extrair as inferências e outras nuances que são dependentes de contexto.

Algumas das técnicas que se classificam como discretas são \textit{one-hot encoding}, \textit{\ac{BoW}}, \textit{\ac{TF-IDF}} e \textit{Count Vectorizer}. \cite{Representacao-Texto:2, Representacao-Texto:3}

As abordagens de representação de texto distribuídas atuam no sentido de representar os elementos componentes de um texto como vetores de números reais com alta densidade. Nesses vetores, cada dimensão corresponde a uma característica específica, embora abstrata, do elemento representado. \cite{Representacao-Distribuida:1} Essas características, as quais são denominadas de \textit{embeddings}, podem não ser diretamente interpretáveis de forma isolada, mas contribuem coletivamente para a codificação de propriedades linguísticas complexas. A adoção de vetores de \textit{embeddings} distribuídos em um espaço vetorial permite que as relações semânticas e sintáticas entre os elementos representados sejam ``espacialmente'' codificadas, de forma que a posição de um vetor em relação a outro, nesse espaço vetorial, captura as relações entre eles. \cite{Representacao-Distribuida:2} Destarte, palavras como ``felino'', ``gato'' e ``bichano'', teriam vetores semelhantes -- e, portanto, distribuídos próximos uns dos outros no espaço vetorial -- posto que compartilham muitas características semânticas. Já palavras como ``gato'' e ``astrolábio'' teriam vetores mais distantes, dado que suas características são menos relacionadas.

O processo de aprendizado dessas representações, codificando o conteúdo textual na forma de \textit{embeddings}, geralmente envolve a análise de grandes conjuntos textuais realizada por modelos, os quais derivam relações contextuais significativas entre os elementos do texto. Destacam-se, entre os modelos mais utilizados para inferir essas representações, o \textit{Word2Vec} \cite{Representacao-Distribuida:3}, o \textit{GloVe} \cite{Representacao-Distribuida:5} e modelos contextuais, como \ac{BERT} e \ac{GPT}. O \textit{Word2Vec} emprega arquiteturas de redes neurais para aprender representações baseadas em contexto, utilizando as abordagens CBOW e \textit{Skip-gram}. O \textit{GloVe}, por sua vez, explora a matriz de coocorrência de palavras para capturar relações semânticas globais. Modelos mais recentes, tais quais \ac{BERT} \cite{Representacao-Distribuida:6} e \ac{GPT} \cite{Representacao-Distribuida:7}, baseados na arquitetura \textit{transformers}, são capazes de aprender representações contextuais mais profundas, capturando nuances de significado dependentes do contexto. Essas representações contextuais são obtidas através de técnicas de pré-treinamento em grandes volumes de texto, permitindo a generalização para diversas tarefas.

Esse conceito de representação distribuída se opõe às representações discretas, sendo mais compacta e eficiente, colaborando para o aprimoramento da viabilidade computacional e o desempenho dos modelos.


\section{Similaridade Textual}
\label{sec:Teorico.3}

Concernente à comparação de textos, com o intuito de mensurar o quanto eles têm em comum, é possível levar em consideração o conteúdo em si -- as mesmas palavras, na mesma estrutura e ordem -- puro e simples, como no caso de detecções de plágio, ou ainda a similaridade do que vai além das palavras, com foco no significado, na intenção do que está escrito, a similaridade semântica. Refere-se, a semelhança semântica, portanto, ao grau de justaposição ou similitude no que tange a significado, entre dois fragmentos de texto, frases, orações ou ainda construtos textuais mais abrangentes, a despeito das palavras -- ou sua ordem -- que os compõem. Observa-se, portanto, que a definição e a avaliação da similaridade semântica envolvem não apenas a identificação de padrões explícitos no texto, mas também a compreensão de aspectos subjacentes relacionados ao significado.

\begin{citacao}
    \textit{Semantic similarity, a subfield within Natural Language Processing (NLP) and computational linguistics, has garnered substantial attention over the past few decades. It operates at the intersection of linguistics and computer science to quantitatively assess the degree of relatedness between linguistic units, ranging from words to sentences and even entire documents. The significance of semantic similarity spans across various domains, including information retrieval, machine translation, text summarization, and recommender systems, among others.} \cite{Similaridade-Semantica:1} \footnote{A similaridade semântica, um subcampo dentro do Processamento de Linguagem Natural (NLP) e da linguística computacional, tem atraído atenção substancial nas últimas décadas. Ela atua na interseção entre a linguística e a ciência da computação para avaliar quantitativamente o grau de relação entre unidades linguísticas, que podem variar de palavras a frases e até documentos inteiros. A importância da similaridade semântica se estende por diversos domínios, incluindo recuperação de informações, tradução automática, sumarização de textos e sistemas de recomendação, entre outros. (Tradução Nossa)}
\end{citacao}


Essa tarefa complexa tem exigido o desenvolvimento de abordagens computacionais cada vez mais sofisticadas, capazes de lidar com a multiplicidade de fatores que influenciam a interpretação semântica. Assim, a transição do uso de métodos tradicionais para técnicas baseadas em abordagens com maior capacidade de apreensão de nuances e contexto, como a utilização de \textit{embeddings} é algo esperado e desejado na área de \ac{PLN}.

\begin{citacao}
    As propostas mais tradicionais para detectar similaridade semântica de textos baseiam-se em análise estrutural, métodos estatísticos e ontológicos. No entanto, como tem acontecido com diversas tarefas de \ac{PLN}, há uma tendência recente de migração para modelos neurais profundos. \cite{Similaridade-Semantica:2}
\end{citacao}

Representar palavras, frases ou documentos como vetores em um espaço de alta dimensionalidade por meio de \textit{embeddings} facilita a quantificação dessa similaridade semântica. A mensuração de similaridade semântica baseada em \textit{embeddings} implica que palavras com significados semelhantes serão representadas por vetores proximamente localizados em sua distribuição no espaço vetorial. A distância entre dois vetores de \textit{embeddings} pode ser calculada por diversas métricas, cada uma delas com suas peculiaridades.

A Distância Euclidiana constitui uma métrica fundamental amplamente empregada em análises que envolvem espaços vetoriais densos, representando a distância linear entre dois pontos em um mesmo (hiper)plano dentro de um espaço \textit{n}-dimensional. Sua utilização destaca-se em cenários em que a magnitude dos vetores possui relevância intrínseca para a interpretação dos dados ou para a modelagem de relações espaciais. Devido às suas propriedades geométricas, a Distância Euclidiana revela-se particularmente apropriada para aplicações em que diferenças absolutas entre vetores, incluindo sua escala, desempenham um papel significativo na análise ou na solução do problema.

Paralelamente, a Distância de Cosseno, diferentemente da métrica Euclidiana, avalia a similaridade entre dois vetores com base no ângulo formado entre eles, sendo particularmente eficaz para capturar a relação de orientação. Essa métrica destaca-se em cenários nos quais a magnitude dos vetores não é tão relevante, residindo o interesse mais fortemente na direção dos vetores. Por essa razão, a Distância de Cosseno é amplamente preferida em diversas aplicações, especialmente naquelas em que a análise da similaridade angular oferece uma representação mais significativa das relações subjacentes entre os vetores.

\section{Bancos Vetoriais (\textit{Vector Stores})}
\label{sec:Teorico.4}
Dada a relevância do papel que a representação de dados textuais via \textit{embeddings} tem para o campo do Processamento de Linguagem Natural, é uma consequência lógica o questionamento de como tornar o acesso a esses dados codificados mais dinâmico e direto, viabilizando agilidade quando necessária a sua utilização. Dessa forma,
\begin{citacao}
    \textit{One of the most common ways to store and search over unstructured data is to embed it and store the resulting embedding vectors, and then at query time to embed the unstructured query and retrieve the embedding vectors that are ``most similar'' to the embedded query.} \cite{Bancos-Vetoriais:1} \footnote{``Uma das formas mais comuns de armazenar e buscar dados não estruturados é criar seus \textit{embeddings} e armazenar os vetores de \textit{embeddings} resultantes. Em seguida, no momento da consulta, também são gerados os \textit{embeddings} da consulta não estruturada, e os vetores de \textit{embeddings} `mais similares' ao da consulta são recuperados''. (Tradução nossa)}
\end{citacao}

A ferramenta utilizada para o armazenamento e busca eficazes de conjuntos de dados textuais e suas representações vetoriais são os \textit{vector stores}. O termo \textit{vector store} designa  um sistema ou plataforma projetada especificamente para lidar com as complexidades e particularidades dos dados vetoriais, frequentemente associada a um \textit{vector database} -- os quais são abordados mais adiante. Esses sistemas são amplamente utilizados em aplicações de aprendizado de máquina, inteligência artificial e análise de dados. Sua capacidade de processar grandes volumes de informações com extrema rapidez os torna ferramentas altamente eficientes para a identificação, organização e recomendação de conteúdos com base em similaridades, especialmente quando integrados a \textit{vector databases}. \cite{Bancos-Vetoriais:2}

Para fins de clareza, convém definir os termos \textit{vector store} e \textit{vector database}, os quais são utilizados frequentemente como conceitos intercambiáveis, mas têm suas diferenças sutis. Enquanto o \textit{vector store} se define nos termos do que já foi apresentado,
\begin{citacao}
    \textit{A vector database, on the other hand, is a more comprehensive system that incorporates the capabilities of a vector store while providing additional features and functionality. (...) Think of them as interconnected components within a larger system. Most available vector storage and retrieval systems are actually vector databases that contain a vector store as a core component.} \cite{Bancos-Vetoriais:3} \footnote{``Um banco de dados vetorial, por outro lado, é um sistema mais abrangente que incorpora as capacidades de um armazenamento vetorial, ao mesmo tempo que oferece recursos e funcionalidades adicionais. (...) Pense neles como componentes interconectados dentro de um sistema maior. A maioria dos sistemas disponíveis de armazenamento e recuperação de vetores são, na verdade, bancos de dados vetoriais que contêm um armazenamento vetorial como componente central''. (Tradução nossa)} 
\end{citacao}

Assim, um \textit{vector database} se trata de uma ferramenta de banco de dados, com funções diversas, o qual compreende um subsistema especializado (o \textit{vector store}) especialmente desenhado para armazenar e buscar dados vetoriais de maneira eficiente. Nesse modelo, o \textit{vector database} atua como uma camada externa para o \textit{vector store}, oferecendo funcionalidades adicionais e ampliando as capacidades de integração com outros sistemas. Ademais, a camada do \textit{vector database} incorpora mecanismos que conectam as operações tradicionais de consulta às funções de recuperação específicas do \textit{vector store}, estabelecendo uma ponte entre as operações clássicas de banco de dados e as demandas de processamento voltadas para dados vetoriais.

Isto posto, neste estudo convencionamos tratar a ambos, \textit{vector store} e \textit{vector database} como banco vetorial ou banco de vetores, dado que, mesmo sendo diferentes, quando se utiliza um \textit{vector database} para realizar operações sobre dados vetoriais, em última instância o que é utilizado é o \textit{vector store} que lhe é interno.

Os bancos vetoriais, portanto, aumentam o aproveitamento de todo o processo de armazenamento e, principalmente, de consulta/recuperação de registros através de busca por similaridade. Enquanto bancos de dados tradicionais, em geral, carecem de eficiência ao lidar com dados vetoriais, os bancos de vetores são otimizados desde o desenho para garantir o bom desempenho do armazenamento e recuperação de vetores de alta dimensionalidade. Para alcançar esse nível de eficácia em dar suporte a dados numéricos multidimensionais e otimização da busca por similaridade, os \textit{vector stores} sacrificam a versatilidade dos bancos convencionais, possuindo esquemas de dados menos flexíveis. \cite{Bancos-Vetoriais:3}

Em seu uso mais regular, ao se criar uma coleção de documentos em um banco vetorial, são definidos tanto uma função de geração de \textit{embeddings} a ser aplicada aos dados inseridos, gerando os vetores a serem armazenados, quanto uma métrica a ser utilizada nas consultas que posteriormente serão feitas. Quando da realização de uma consulta, portanto, os dados de consulta utilizados (a \textit{query}) são processados pela mesma função de geração de \textit{embeddings} utilizada para gerar a representação das instâncias armazenadas, sendo sua representação vetorial utilizada para recuperar os registros mais semelhantes, de acordo com a métrica previamente estabelecida.





\section{\textit{Large Language Models} (LLMs)}
\label{sec:Teorico.5}

\textit{Large Language Models} ou LLMs (ocasionalmente chamados de Modelos de Linguagem de Grande Escala), são modelos estatísticos avançados, gerados por aprendizado de máquina, os quais são projetados para compreender e gerar linguagem humana. Eles realizam essa tarefa por meio da identificação de padrões e relações complexas entre palavras. Sua funcionalidade central está baseada na previsão da próxima palavra ou fragmento de palavra (\textit{token}) em uma sequência, considerando o contexto fornecido pelos \textit{tokens} anteriores. Para atingir esse objetivo, os \ac{LLM}s são treinados em grandes volumes de dados textuais, o que lhes permite identificar padrões e estruturas complexas na linguagem natural. \cite{LLMs:1, LLMs:2}

O desenvolvimento de \ac{LLM}s envolve processos intensivos de treinamento baseados em técnicas de aprendizado de máquina, especialmente o aprendizado não supervisionado, que trata de aprender padrões e semelhanças entre elementos diversos de um conjunto de dados, sem intervenção externa. Nesse método, os modelos aprendem os significados das palavras a partir de suas coocorrências em grandes conjuntos de dados textuais. Durante o treinamento, os modelos realizam operações estatísticas sobre representações numéricas dos textos, no formato de vetores de \textit{embeddings}, que representam palavras e suas características de distribuição, realizando o ajuste de seus parâmetros para minimizar os erros nas previsões de próximos \textit{tokens}. Dessa forma, por meio da avaliação da probabilidade de ocorrência de uma sequência de palavras/\textit{tokens}, são mantidas as informações de contexto e coerência entre as partes do texto. É precisamente por isso que um dos componentes fundamentais no funcionamento dos \ac{LLM}s é o mecanismo de atenção (\textit{attention mechanism}), que avalia a relevância de \textit{tokens} anteriores em relação ao contexto atual, em favor de respostas coerentes e contextualmente apropriadas.

Nesse contexto, o desenvolvimento de técnicas específicas e de \textit{embeddings} pré-treinados, como GloVe e fastText, trouxe avanços significativos na representação semântica das palavras, permitindo a aplicação de técnicas de aprendizado por transferência (\textit{transfer learning}). Esses processos possibilitam que os \ac{LLM}s  compreendam a linguagem de forma contextualizada e escalável, tornando-os altamente eficientes em tarefas linguísticas complexas, fazendo dos \ac{LLM}s ferramentas versáteis em tarefas de processamento de linguagem natural, como classificação de texto e previsão contextual.

Os \ac{LLM}s possuem uma ampla e diversificada gama de aplicações práticas, que incluem a análise de sentimento de um texto (positivo ou negativo), a predição de palavras mais prováveis a seguir em um contexto específico, a geração de respostas adequadas para perguntas, geração de texto com base em uma descrição do que se espera que seja a saída e uma miríade de outras tarefas. Por exibirem essas capacidades, tais modelos têm sido amplamente utilizados em ferramentas como \textit{chatbots}, sistemas de geração de conteúdo e outras soluções baseadas em inteligência artificial. Sua capacidade de lidar com tarefas variadas e produzir textos com fluidez torna-os indispensáveis em contextos como o da pesquisa acadêmica e estudos de comunicação política, dado que permitem a análise e classificação textual, ampliando a compreensão sobre o uso da linguagem nesse contexto. Também contribuem para superar desafios relacionados à reprodutibilidade e validade de pesquisas, gerando textos semelhantes aos produzidos por humanos. Isso é especialmente relevante em tarefas como criação de conteúdo e respostas automatizadas, onde modelos como o \ac{GPT}, popularizado com a ferramenta ChatGPT, têm ganhado destaque público, consolidando-se como expoentes da inteligência artificial contemporânea \cite{LLMs:3,LLMs:4}.

Apesar de seus avanços, os \ac{LLM}s  enfrentam desafios significativos. Um dos principais é a presença de vieses nos dados de treinamento, o que pode levar à perpetuação de preconceitos e discriminações. Considerações éticas também emergem na aplicação dessas tecnologias, incluindo a necessidade de validação rigorosa para garantir que suas interpretações sejam consistentes com a compreensão humana. A famosa máxima de que ``todos os modelos quantitativos de linguagem estão errados — mas alguns são úteis'' ilustra as limitações inerentes dos \ac{LLM}s, reforçando a importância de uma avaliação crítica de seus resultados.

Um equívoco comum é a crença de que os \ac{LLM}s são meros preditores de próximo \textit{token}. Embora a previsão do próximo termo de fato seja o resultado final de seu funcionamento, essa visão simplista subestima suas capacidades \cite{LLMs:2}. De forma a realizar a predição de \textit{tokens} individuais, que permite a geração de textos estruturados e significativos, os \ac{LLM}s aprendem associações complexas entre \textit{tokens} (os quais, em conjunto, compõem conceitos), formando redes de significados e interrelações, as quais são ``persistidos'' intrinsecamente em sua estrutura interna. Essa sofisticação vai além de meras previsões estatísticas, envolvendo a codificação de informações estruturais mais elevadas, como aspectos sintáticos e semânticos, mapeados em um conjunto de significados e relações entre conceitos.

Conforme posto por \citeonline{LLMs:2}:
\begin{citacao}
    \textit{The claim that LLMs are just next token predictors blinds us to the capabilities of reinforcement learning to achieve this benchmark, and undermines our ability to explain how such structural information is learned by the model. A sign that this higher level structural information -- e.g. inter- and intra- sentential syntactic information -- is present in LLMs is revealed by the expansive attention heads the models use to make next token predictions. Predicting the next token, conditional on the previous token only, is the simplest form of prediction. If LLMs did only this, they would not be very successful at answering questions or telling stories and would perhaps indeed be mere next token predictors.} \footnote{A afirmação de que os LLMs são apenas preditores de próximo \textit{token} nos impede de perceber as capacidades do aprendizado por reforço em alcançar esse marco e prejudica nossa capacidade de explicar como essas informações estruturais são aprendidas pelo modelo. Um sinal de que essas informações estruturais de nível mais alto -- como informações sintáticas inter e intra-sentenciais -- estão presentes nos LLMs é revelado pelos \textit{attention headers} expansivos que os modelos utilizam para fazer previsões de próximos \textit{tokens}. Prever o próximo \textit{token}, condicionado apenas ao \textit{token} anterior, é a forma mais simples de previsão. Se os LLMs fizessem apenas isso, eles não seriam muito bem-sucedidos em responder perguntas ou contar histórias e, talvez, de fato fossem meros preditores de próximo \textit{token}. (Tradução nossa)}
\end{citacao}


\section{\textit{Retrieval-Augmented Generation} (RAG)}
\label{sec:Teorico.6}

Os \ac{LLM}s, assim como outros modelos baseados em redes neurais, são capazes de aprender e armazenar conhecimento na forma de interconexões e relacionamentos entre os dados, sendo essas informações armazenadas em seu próprio conjunto de parâmetros. Após o treinamento, esses modelos tornam-se independentes de fontes externas, funcionando como uma base de conhecimento autocontida. Embora essa abordagem seja poderosa, ela apresenta algumas limitações consideráveis.

Dado que o conhecimento apreendido é armazenado diretamente na estrutura do modelo, atualizá-lo ou expandi-lo exige modificações significativas. Além disso, os mecanismos que sustentam suas predições são complexos e difíceis de interpretar, o que pode resultar em saídas imprecisas ou equivocadas, comumente chamadas de ``alucinações''. Assim, a despeito de as poderosas capacidades dos \ac{LLM}s os tornarem especialmente úteis para responder rapidamente a solicitações gerais, sua funcionalidade apresenta restrições, em especial quando se trata de realizar uma análise mais detalhada sobre tópicos específicos ou atualizados, fora do escopo de seu treinamento.

Havendo apontado esses pontos negativos a serem contornados, \citeonline{RAG:1} sugerem que modelos híbridos, que combinam memória armazenada em parâmetros (como os \ac{LLM}s) com conhecimento externo, obtido por meio de recuperação de dados vindos de uma base de conhecimento, podem solucionar algumas dessas questões. Apontam que, dessa forma, ``o conhecimento pode ser diretamente revisado e expandido, podendo o conteúdo acessado ser inspecionado e interpretado''. \footnote{``\textit{knowledge can be directly revised and expanded, and accessed knowledge can be inspected and interpreted}'' (Tradução nossa)}.

A técnica proposta por \citeonline{RAG:1}, a qual foi nomeada como \textit{\ac{RAG}} (Geração aumentada por recuperação) consiste precisamente, conforme sugerido no próprio nome, em aprimorar a geração de respostas do \ac{LLM} por meio de incluir no \textit{prompt} que a ele é enviado informações que o ajudem a formular sua saída com base em dados corretos.

Para isso, é adicionada uma camada no processo de envio do \textit{input} do usuário ao \ac{LLM}, que se encarrega de realizar o seguinte:

\begin{itemize}
    \item utiliza o \textit{input} para realizar uma busca de similaridade semântica em um banco vetorial, previamente populado com fragmentos de documentos pertinentes ao domínio que a aplicação se propõe a abordar;
    \item gera um \textit{prompt} que inclui diretrizes para o \ac{LLM}, orientando-o a utilizar os resultados do banco de vetores como corpo de conhecimento para gerar a resposta;
    \item encaminha o \textit{prompt} aumentado ao \ac{LLM}.
\end{itemize}

A incorporação desse conhecimento obtido de fontes externas confiáveis no processo de geração de resposta reduz drasticamente a probabilidade de geração de conteúdo fatalmente incorreto por parte dos \ac{LLM}s. Não somente aumenta a confiabilidade, outrossim amplia a possibilidade de expansão e atualização das informações disponíveis para a geração das respostas, uma vez que a inclusão de novas informações nas bases externas de conhecimento, assim como a correção de eventuais informações incorretas ou obsoletas, se demonstra um procedimento menos custoso, menos demorado e mais flexível.

Tais características promissoras justificam o que foi apontado por \citeonline{RAG:2}, ao afirmar que:

\begin{citacao}
    \textit{RAG technology has rapidly developed in recent years (...) Its integration into LLMs has resulted in widespread adoption, establishing RAG as a key technology in advancing chatbots and enhancing the suitability of LLMs for real-world applications.} \footnote{A tecnologia RAG se desenvolveu rapidamente nos últimos anos (...) Sua integração aos LLMs resultou em uma ampla adoção, estabelecendo a RAG como uma tecnologia chave para o avanço de chatbots e para aumentar a adequação dos LLMs em aplicações do mundo real. (Tradução nossa)}
\end{citacao}

\section{Métricas de Avaliação do Texto Gerado}
\label{sec:Teorico.7}

A avaliação da qualidade do texto gerado por \ac{LLM}s pode ser levada a cabo inteiramente por analistas humanos, todavia a empreitada como um todo seria cara em termos de tempo despendido e custo financeiro, não sendo, logo, escalável e sustentável. Por isto, sua complementação com a avaliação feita por meios automatizados é uma prática comum. \cite{MetricasAvaliacaoTexto:1} Dentre o conjunto de ferramentas utilizadas para medição da qualidade do texto sintético, as métricas baseadas em referência se colocam em um lugar de relevância na avaliação de modelos de \ac{PLN}, especialmente em tarefas que envolvem geração e tradução de texto. Essas métricas avaliam a qualidade do texto gerado comparando-o diretamente com referências anotadas por humanos.

Entre elas, destacam-se as métricas baseadas em n-gramas, como BLEU (\textit{Bilingual Evaluation Understudy}) e ROUGE (\textit{Recall-Oriented Understudy for Gisting Evaluation}). Ao analisar sequências sobrepostas de palavras ou unidades sublexicais (n-gramas), essas métricas fornecem recursos para análises quantitativas sobre a fidelidade linguística e o alinhamento semântico das saídas geradas pelos modelos em relação aos textos de referência. Apesar de terem sido desenvolvidas em paradigmas anteriores do \ac{PLN}, sua aplicabilidade aos \ac{LLM}s evidencia sua relevância duradoura.

A pontuação BLEU, um marco consolidado na avaliação de tradução automática, quantifica a precisão do texto gerado avaliando a sobreposição de n-gramas com traduções de referência, fornecidas/notadas por humanos. \cite{BLEU:1} Inicialmente concebido para avaliar traduções bilíngues, o BLEU foi ampliado para diversos domínios, incluindo a sumarização e a geração de paráfrases. Sua aplicação se baseia no cálculo da proporção de n-gramas no texto candidato que aparecem no texto de referência, sendo utilizada a média calculada dos resultados ao longo de um corpus para gerar uma pontuação de qualidade geral. Notavelmente, o BLEU incorpora n-gramas de ordem superior (por exemplo, bigramas ou trigramas) para capturar a precisão contextual, enquanto atenua a dependência excessiva de unigramas. Embora seja raro que traduções humanas atinjam pontuações perfeitas, a capacidade da métrica de facilitar comparações entre sistemas ressalta sua utilidade na padronização de práticas de avaliação em diversas aplicações.

O ROUGE, por outro lado, é uma métrica focada no \textit{recall}, que enfatiza a completude das informações capturadas no texto gerado em relação ao texto de referência. \cite{ROUGE:1} Empregado predominantemente em tarefas de sumarização, o ROUGE mede o grau em que o conteúdo essencial é retido. O ROUGE calcula precisão, \textit{recall} e \textit{f1-score} do texto candidato em relação ao texto de referência. Considerando que o ROUGE é orientado ao \textit{recall}, se demonstra, ele, particularmente adequado para aplicações em que a captura de informações abrangentes é crucial, complementando métricas orientadas à precisão, como o BLEU, e fornecendo uma perspectiva holística sobre a qualidade textual.

Todavia, as métricas tradicionais baseadas em n-grams, como BLEU e ROUGE, enfrentam desafios consideráveis ao lidar com textos que apresentam variação significativa em relação às referências. Uma de suas limitações se refere à identificação de paráfrases, conforme apontado por \citeonline{BERTScore:1}, ao apontar que:

\begin{citacao}
    \textit{(...) such methods often fail to robustly match paraphrases. For example, given the reference \textit{people like foreign cars}, BLEU and METEOR incorrectly give a higher score to \textit{people like visiting places abroad} compared to \textit{consumers prefer imported cars}. This leads to performance underestimation when semantically-correct phrases are penalized because they differ from the surface form of the reference.}
    \footnote{[...] tais métodos frequentemente falham em correlacionar paráfrases de forma robusta. Por exemplo, dada a referência \textit{as pessoas gostam de carros do exterior}, o BLEU e o METEOR incorretamente atribuem uma pontuação mais alta para \textit{as pessoas gostam de visitar lugares no exterior} em comparação com \textit{consumidores preferem carros importados}. Isso leva ao sub-desempenho quando frases semanticamente corretas são penalizadas por diferirem da forma superficial da referência.}
\end{citacao}

Além disso, por se basear em n-gramas, tais métricas não conseguem capturar relações de dependência semântica que não estão tão próximas no texto, se estendendo para além do tamanho do n-grama, bem como se demonstram incapazes de perceber adequadamente alterações significativas na ordem semântica. Por exemplo, uma troca de ``A porque B'' para ``B porque A'' pode não impactar substancialmente a pontuação BLEU, mesmo que a mudança semântica seja relevante. 

Como solução para essas limitações, o BERTScore foi desenvolvido, introduzindo uma abordagem baseada em embeddings contextuais, se propondo a ser

\begin{citacao}
    \textit{a new metric for evaluating generated text against gold standard references. BERTSCORE is purposely designed to be simple, task-agnostic, and easy to use. Our analysis illustrates how BERTSCORE resolves some of the limitations of commonly used metrics, especially on challenging adversarial examples.} \cite{BERTScore:1}
    \footnote{uma nova métrica para a avaliação de textos gerados em comparação com textos de referência. O BERTSCORE foi intencionalmente projetado para ser simples, independente da natureza da tarefa e fácil de usar. Nossa análise ilustra como o BERTSCORE resolve algumas das limitações das métricas comumente utilizadas, especialmente em exemplos adversariais desafiadores. (Tradução nossa)}
\end{citacao}

A métrica calcula a similaridade entre dois textos, sendo um a referência anotada, \(x\), e o outro um fragmento de texto gerado, \(\hat{x}\), utilizando as similaridades de cosseno entre suas representações por \textit{embeddings}. O \textit{score} completo faz uma comparação dos \textit{embeddings} de cada \textit{token} em \(x\) com os \textit{embeddings} de cada \textit{token} em \(\hat{x}\), buscando pelo par de \textit{tokens} com maior similaridade de cosseno para calcular o recall. A precisão, similarmente, é calculada ao se realizar uma busca pelos tokens de  presentes em x, para calcular a precisão. O F1 score é calculado pela combinação da precisão e recall calculados.

\[
R_{\text{BERT}} = \frac{1}{|x|} \sum_{x_i \in x} \max_{\hat{x}_j \in \hat{x}} \text{x}_i^{\top} \hat{\text{x}}_j
\]

\[
P_{\text{BERT}} = \frac{1}{|\hat{x}|} \sum_{\hat{x}_j \in \hat{x}} \max_{x_i \in x} \text{x}_i^{\top} \hat{\text{x}}_j
\]

\[
F_{\text{BERT}} = 2 \frac{P_{\text{BERT}} \cdot R_{\text{BERT}}}{P_{\text{BERT}} + R_{\text{BERT}}}
\]

Essa abordagem permite uma avaliação mais robusta, uma vez que \textit{embeddings} pré-treinados capturam nuances contextuais e semânticas das palavras em suas respectivas frases. O BERTScore, portanto, é capaz de reconhecer paráfrases e lidar com dependências semânticas distantes, superando limitações que comprometem a eficácia de métricas tradicionais.

Semelhantemente ao BLEU e ao ROUGE, o BERTScore fundamenta-se em duas métricas principais: precisão e \textit{recall}. A precisão avalia o quanto o conteúdo do texto candidato é relevante em relação ao texto de referência, de forma que uma alta precisão indica que a maior parte da informação no texto candidato está semanticamente alinhada com o texto de referência, refletindo a fidelidade do conteúdo gerado. O \textit{recall}, por sua vez, mensura o quanto do conteúdo presente no texto de referência é capturado pelo texto candidato; assim, um valor alto de \textit{recall} sugere que o texto gerado cobre bem as informações essenciais do texto de referência, evidenciando sua completude.

Combinando precisão e abrangência, o BERTScore oferece uma métrica sofisticada que, além de avaliar similaridade, captura as nuances semânticas necessárias para tarefas como tradução, sumarização e parafraseamento. Dessa forma, destaca-se como uma alternativa avançada às métricas baseadas em n-gramas, particularmente em cenários que exigem uma análise mais refinada e sensível ao contexto linguístico.

